\section{Discussion}
\label{sec:discussion-and-limits}

\paragraph{What the results say.}
In the setting we studied, a model can tell from its own answers what a round
of self-improvement will do, select its own training data, and decide where to
train. Two quantities carry all three roles: \emph{novelty} (how often the
teacher's answers differ from the student's) and the \emph{confirmed share}
(how often a second split reproduces them), standing in for the correct answers
and hand-written filter used by \citet{shin2026compositional}. As a filter,
agreement between splits is what lets the loop work at all: without it, the
model learns nothing from direct answers or learns the composition's carry
error. As a measurement, novelty tracks how much of each length remains to be
learned closely enough to skip a third to a half of the work at no cost in
accuracy---provided any cutoff placed on it matches the question asked
(flagging a large single-round jump versus checking whether any gap remains).
We stopped the ladder at 39 digits by choice: on two starting models, final
accuracy at the longest sums matched that at the shortest, while on the third
errors accumulated slowly, leaving open how much further it would have climbed.

\paragraph{Relation to other work.}
\label{sec:related}
Training a model on outputs produced with an auxiliary procedure---tree search
\citep{anthony2017thinking}, decomposition \citep{christiano2018amplification},
or richer prompts \citep{zhang2026basersi}---is standard practice, though
training on self-generated outputs can also degrade a model
\citep{shumailov2024collapse}. Whereas prior work uses the teacher--student gap
to drive training, or predicts self-training gains from dataset statistics
\citep{vanasch2016predicting}, we use that gap before training to decide
whether and where to train. Two concurrent studies ask related questions with
stronger signals: \citet{li2026evolve} select self-distillation data using
gradients and ground-truth labels, and \citet{he2026recursive} filter agent
generations only after training and measuring their gain; we decide beforehand
from outputs alone. For estimating accuracy without labels, prior work relies
on agreement across models \citep{baek2022agreement}, across two views in
co-training \citep{blum1998cotraining}, or among a model's own samples at test
time \citep{wang2023selfconsistency} and during self-training
\citep{shafayat2025selftrain}. Our sample-and-vote baseline follows the last
approach and taught nothing; agreement across distinct splits acts instead like
two views inside a single model, informative precisely because the computation
paths and their errors differ. Finally, on length generalization, removing
positional encodings \citep{kazemnejad2023positional} or adding position-aware
embeddings \citep{mcleish2024abacus} extends addition
architecturally---\citet{zhou2024algorithms} characterize which algorithms
extrapolate in this way---while filtered self-training on direct answers
extends our architecture one digit at a time \citep{lee2025selfimproving}; here
the architecture stays fixed and split-confirmed composition carries it across
32 lengths.

\paragraph{Scope and outlook.}
\label{sec:limits}
Our empirical evidence is of two kinds: the single-round predictor was frozen
before evaluation, whereas the 39-digit ladder is exploratory---its rules were
fixed before running, save for the divergence retry rule added on the first
seed, and its pacing threshold was set from logs of an earlier loop. Using one
small model, three evaluation seeds, and a synthetic task was a deliberate
choice so every output could be scored exactly and harmful rounds induced
cleanly. Extending the same two quantities to larger models and open-ended
tasks will test three assumptions built into our setup. First, the checks
evaluate a supplied procedure that admits distinct splits; they do not discover
how to decompose a task. Second, they inspect only the answers a teacher keeps,
not what selection omits: a teacher that is right on every answer it keeps
can still harm the student by leaving a whole class of problems out of
training, and neither quantity would show it. Third, exact
string match will need to give way to semantic agreement where valid answers
are not unique.
